\documentclass[a4paper,oneside]{book}

%----------------------------------------------------------------------------------------
%	1. GÓI NGÔN NGỮ & FONT
%----------------------------------------------------------------------------------------
% \usepackage[utf8]{inputenc}
\usepackage[T5,T1]{fontenc}
\usepackage{newtxtext, newtxmath} % Font Times New Roman
\usepackage[vietnamese]{babel}

\usepackage{booktabs,tabularx}
\usepackage{algorithm}
%----------------------------------------------------------------------------------------
%	2. CÁC GÓI CƠ BẢN
%----------------------------------------------------------------------------------------
\usepackage{amsmath}
\usepackage{graphicx}
\usepackage{array,booktabs,multirow}
\usepackage{siunitx} % canh lề số đẹp
\sisetup{
  round-mode=places,round-precision=3,
  table-number-alignment = center,
  detect-weight=true}

\usepackage{subcaption}
\usepackage{float}
\usepackage{setspace}
\usepackage{longtable}
\usepackage{tikz}
\usetikzlibrary{calc}
\usepackage{etoolbox}

%----------------------------------------------------------------------------------------
%	3. BỐ CỤC & SIÊU LIÊN KẾT
%----------------------------------------------------------------------------------------
\usepackage{geometry}
\geometry{
  a4paper,
  top=20mm,
  bottom=20mm,
  left=30mm,
  right=20mm,
  headheight=18pt,
  headsep=10mm,
  footskip=10mm,
}
\definecolor{cobalt}{rgb}{0.0, 0.28, 0.67}
\definecolor{carmine}{rgb}{0.59, 0.0, 0.09}
\usepackage[unicode, colorlinks=true, urlcolor=carmine, citecolor=cobalt, linkcolor=carmine]{hyperref}

%----------------------------------------------------------------------------------------
%	4. ĐỊNH DẠNG TIÊU ĐỀ (CHƯƠNG, MỤC, CAPTION)
%----------------------------------------------------------------------------------------
\usepackage{titlesec}
\usepackage{caption}
\setcounter{secnumdepth}{3} % đánh số tới \subsubsection  -> 1.3.1
\setcounter{tocdepth}{3}    % hiển thị tới \subsubsection trong TOC

% Định dạng tiêu đề Chương
\titleformat{name=\chapter}[block]
  {\normalfont\fontsize{14pt}{20pt}\bfseries\centering}
  {Chương \thechapter.}{1em}{\MakeUppercase}
\titleformat{name=\chapter,numberless}[block]
  {\normalfont\fontsize{14pt}{20pt}\bfseries\centering}
  {}{0em}{\MakeUppercase}
\titlespacing*{\chapter}{0pt}{-10pt}{14pt}

% Định dạng tiêu đề Section & Subsection
\titleformat{\section}[hang]
  {\normalfont\fontsize{13pt}{15pt}\selectfont\bfseries}
  {\thesection}{1em}{}
\titleformat{\subsection}[hang]
  {\normalfont\fontsize{13pt}{16.5pt}\selectfont\bfseries}
  {\thesubsection}{1em}{}

% Định dạng tiêu đề cho hình, bảng
\captionsetup{font=normalsize, labelfont=bf, labelsep=colon}

%----------------------------------------------------------------------------------------
%	5. ĐỊNH DẠNG MỤC LỤC & CÁC DANH SÁCH
%----------------------------------------------------------------------------------------
\usepackage{tocloft}

% Chuẩn hóa tên các trang danh mục theo đúng thuật ngữ chính thức, ghi đè sau khi
% babel[vietnamese] đã khởi tạo giá trị mặc định của nó (dùng \AtBeginDocument để
% không bị babel ghi đè lại lúc \begin{document}).
\AtBeginDocument{%
  \renewcommand{\contentsname}{Mục lục}
  \renewcommand{\listfigurename}{Danh mục các hình}
  \renewcommand{\listtablename}{Danh mục các bảng}
}

% Định dạng tiêu đề của các trang danh sách: in đậm, IN HOA, cỡ 14
\makeatletter
\renewcommand{\@cftmaketoctitle}{%
  \clearpage
  {\centering\bfseries\fontsize{14pt}{20pt}\selectfont \MakeUppercase{\contentsname}\par}
  \nobreak\vskip 14pt
}
\renewcommand{\@cftmakeloftitle}{%
  \clearpage
  {\centering\bfseries\fontsize{14pt}{20pt}\selectfont \MakeUppercase{\listfigurename}\par}
  \nobreak\vskip 14pt
}
\renewcommand{\@cftmakelottitle}{%
  \clearpage
  {\centering\bfseries\fontsize{14pt}{20pt}\selectfont \MakeUppercase{\listtablename}\par}
  \nobreak\vskip 14pt
}
\makeatother

% Định dạng các dòng trong Mục lục
\renewcommand{\cftchappresnum}{Chương }
\renewcommand{\cftchapaftersnum}{.}
\setlength{\cftchapnumwidth}{8em}
\renewcommand{\cftsecleader}{\cftdotfill{\cftdotsep}}
\renewcommand{\cftsubsecleader}{\cftdotfill{\cftdotsep}}
\renewcommand{\cftdotsep}{3}

% Làm đều khoảng cách các dòng trong Mục lục
\setlength{\cftbeforechapskip}{\parskip}
\setlength{\cftbeforesecskip}{\parskip}
\setlength{\cftbeforesubsecskip}{\parskip}

%----------------------------------------------------------------------------------------
%	6. CÁC GÓI MỞ RỘNG (CODE, GIẢI THUẬT)
%----------------------------------------------------------------------------------------
% (gói algorithm đã được nạp ở mục 1 phía trên, không nạp lại ở đây để tránh trùng lặp)
\usepackage{algpseudocode}
\makeatletter
\renewcommand{\ALG@name}{Giải thuật}
\renewcommand{\listalgorithmname}{Danh sách giải thuật}
\makeatother

\usepackage{listings}
\usepackage{xcolor}
\lstset{
  language=C++,
  basicstyle=\ttfamily\footnotesize,
  keywordstyle=\color{blue}\bfseries,
  stringstyle=\color{red},
  commentstyle=\color{green!50!black}\itshape,
  breaklines=true,
  frame=single,
  numbers=left,
  numberstyle=\tiny\color{gray},
  tabsize=2
}

%----------------------------------------------------------------------------------------
%	7. FONT CHUNG & KHOẢNG CÁCH ĐOẠN
%----------------------------------------------------------------------------------------
% Cỡ chữ 13pt, giãn dòng ~1.15 (13pt x 1.15 ~ 15pt baselineskip) theo yêu cầu
% "line spacing 1.1 to 1.2, use around 1.15". Không dùng thêm \setstretch của gói
% setspace ở đây để tránh nhân đôi hệ số giãn dòng (chỉ dùng setspace cục bộ cho
% các trang mục lục/danh mục nếu cần).
\makeatletter
\renewcommand\normalsize{%
  \@setfontsize\normalsize{13pt}{14.95pt}%
  \abovedisplayskip 6pt plus 2pt minus 2pt
  \abovedisplayshortskip \z@ plus 2pt
  \belowdisplayshortskip 6pt plus 2pt minus 2pt
  \belowdisplayskip \abovedisplayskip
}
\makeatother
\normalsize
% Khoảng cách đoạn trước/sau ~3pt theo yêu cầu định dạng
\setlength{\parskip}{3pt}
\setlength{\parindent}{15pt}

\usepackage{tocloft}

\renewcommand{\cftfigpresnum}{Hình~}
\setlength{\cftfignumwidth}{4.0em}

\renewcommand{\cfttabpresnum}{Bảng~}
\setlength{\cfttabnumwidth}{4.0em}

\usepackage{thesis}
\usepackage{setspace}

\cftsetpnumwidth{2em}
\cftsetrmarg{3em}
\setlength{\cftsecnumwidth}{3em}
\setlength{\cftsubsecnumwidth}{4em}